% GENERATED FILE. Edit canonical lessons, then run npm run generate:books.
% canonical-source-hash: fnv1a64:91b1e1fb56277cc3
% canonical-lessons: SA-C204-sucayati, SA-C204-khanati, SA-C204-vapati, SA-C204-jivati, SA-C204-tapati

\chapter{Points out, Digs, Sows, Is alive, Heats}
\label{ch:sa-points-out-digs-sows-is-alive-heats}
\hlchaptermodality{204}

\begin{chapteropening}
\textbf{By the end of this chapter:} \emph{I can say points out, digs, sows, is alive and heats.}

Complete the last lesson of chapter 204: I can say points out, digs, sows, is alive and heats.
\end{chapteropening}

\section[\texorpdfstring{\sk{सूचयति} (sūcayati)}{सूचयति (sūcayati)}]{\sk{सूचयति} (sūcayati) — points out, informs}

\label{lesson:SA-C204-sucayati}



\begin{quote}
Before the new one: say the Sanskrit for speaks, then the Sanskrit for shines.
\end{quote}

\subsection*{You'll want to know: \sk{सूचयति}}
\textbf{\sk{सूचयति}} — \emph{sūcayati} — \textquotedblleft{}points out, informs\textquotedblright{}.

\begin{grammarlens}[title={What you've built}]
One more word for this chapter.
\end{grammarlens}

\subsection*{Your turn}
\begin{itemize}
\raggedright
  \item \emph{Say it:} \emph{sūcayati}
  \item \emph{Say it:} \emph{sūcayati}, once more
  \item \emph{Say it:} say \emph{sūcayati}
  \item \emph{Recall:} say the Sanskrit for speaks, then the Sanskrit for shines, then say \emph{sūcayati} again
\end{itemize}

\begin{tcolorbox}[breakable,colback=teal!4,colframe=teal!35!black,title={Before you move on}]
What does \sk{सूचयति} mean? (\textquotedblleft{}Points out, informs\textquotedblright{}.) Say it once more.
\end{tcolorbox}
\section[\texorpdfstring{\sk{खनति} (khanati)}{खनति (khanati)}]{\sk{खनति} (khanati) — digs}

\label{lesson:SA-C204-khanati}



\begin{quote}
Before the new one: say the Sanskrit for shines, then the Sanskrit for points out.
\end{quote}

\subsection*{You'll want to know: \sk{खनति}}
\textbf{\sk{खनति}} — \emph{khanati} — \textquotedblleft{}digs\textquotedblright{}.

\begin{grammarlens}[title={What you've built}]
One more word for this chapter.
\end{grammarlens}

\subsection*{Your turn}
\begin{itemize}
\raggedright
  \item \emph{Say it:} \emph{khanati}
  \item \emph{Say it:} \emph{khanati}, once more
  \item \emph{Say it:} say \emph{khanati}
  \item \emph{Recall:} say the Sanskrit for shines, then the Sanskrit for points out, then say \emph{khanati} again
\end{itemize}

\begin{tcolorbox}[breakable,colback=teal!4,colframe=teal!35!black,title={Before you move on}]
What does \sk{खनति} mean? (\textquotedblleft{}Digs\textquotedblright{}.) Say it once more.
\end{tcolorbox}
\section[\texorpdfstring{\sk{वपति} (vapati)}{वपति (vapati)}]{\sk{वपति} (vapati) — sows}

\label{lesson:SA-C204-vapati}



\begin{quote}
Before the new one: say the Sanskrit for points out, then the Sanskrit for digs.
\end{quote}

\subsection*{You'll want to know: \sk{वपति}}
\textbf{\sk{वपति}} — \emph{vapati} — \textquotedblleft{}sows\textquotedblright{}.

\begin{grammarlens}[title={What you've built}]
One more word for this chapter.
\end{grammarlens}

\subsection*{Your turn}
\begin{itemize}
\raggedright
  \item \emph{Say it:} \emph{vapati}
  \item \emph{Say it:} \emph{vapati}, once more
  \item \emph{Say it:} say \emph{vapati}
  \item \emph{Recall:} say the Sanskrit for points out, then the Sanskrit for digs, then say \emph{vapati} again
\end{itemize}

\begin{tcolorbox}[breakable,colback=teal!4,colframe=teal!35!black,title={Before you move on}]
What does \sk{वपति} mean? (\textquotedblleft{}Sows\textquotedblright{}.) Say it once more.
\end{tcolorbox}
\section[\texorpdfstring{\sk{जीवति} (jīvati)}{जीवति (jīvati)}]{\sk{जीवति} (jīvati) — is alive, lives}

\label{lesson:SA-C204-jivati}



\begin{quote}
Before the new one: say the Sanskrit for digs, then the Sanskrit for sows.
\end{quote}

\subsection*{You'll want to know: \sk{जीवति}}
\textbf{\sk{जीवति}} — \emph{jīvati} — \textquotedblleft{}is alive, lives\textquotedblright{}.

\begin{grammarlens}[title={What you've built}]
One more word for this chapter.
\end{grammarlens}

\subsection*{Your turn}
\begin{itemize}
\raggedright
  \item \emph{Say it:} \emph{jīvati}
  \item \emph{Say it:} \emph{jīvati}, once more
  \item \emph{Say it:} say \emph{jīvati}
  \item \emph{Recall:} say the Sanskrit for digs, then the Sanskrit for sows, then say \emph{jīvati} again
\end{itemize}

\begin{tcolorbox}[breakable,colback=teal!4,colframe=teal!35!black,title={Before you move on}]
What does \sk{जीवति} mean? (\textquotedblleft{}Is alive, lives\textquotedblright{}.) Say it once more.
\end{tcolorbox}
\section[\texorpdfstring{\sk{तपति} (tapati)}{तपति (tapati)}]{\sk{तपति} (tapati) — heats, is hot}

\label{lesson:SA-C204-tapati}



\begin{quote}
Before the new one: say the Sanskrit for sows, then the Sanskrit for is alive.
\end{quote}

\subsection*{You'll want to know: \sk{तपति}}
\textbf{\sk{तपति}} — \emph{tapati} — \textquotedblleft{}heats, is hot\textquotedblright{}.

\begin{grammarlens}[title={What you've built}]
One more word for this chapter.
\end{grammarlens}

\subsection*{Your turn}
\begin{itemize}
\raggedright
  \item \emph{Say it:} \emph{tapati}
  \item \emph{Say it:} \emph{tapati}, once more
  \item \emph{Say it:} say \emph{tapati}
  \item \emph{Recall:} say the Sanskrit for sows, then the Sanskrit for is alive, then say \emph{tapati} again
\end{itemize}

\begin{tcolorbox}[breakable,colback=teal!4,colframe=teal!35!black,title={Before you move on}]
What does \sk{तपति} mean? (\textquotedblleft{}Heats, is hot\textquotedblright{}.) Say it once more.
\end{tcolorbox}
